\subsection{Sufficiency, Exponential Families, and Local Asymptotics}

The shape of a model determines which summaries of the data preserve information and which approximations become valid in large samples.

\begin{definition} \label{def:statistics-sufficiency-ancillarity} \hlt{Statistics, Sufficiency, and Ancillarity}
\begin{enumerate}[label=(\roman*)]
\setlength{\itemsep}{0pt}
\item A \hlt{statistic} is a measurable map $T:\mathcal{X}\to\mathcal{T}$.
\item $T$ is \hlt{sufficient} if $\mathcal{L}_\theta(X\mid T)$ is independent of $\theta$.
\item $T$ is \hlt{ancillary} if $\mathcal{L}_\theta(T)$ is independent of $\theta$.
\end{enumerate}
\end{definition}

\begin{theorem} \label{thm:factorization-criterion-sufficiency} \hlt{Factorization Criterion}\\
Let $\mathcal{P}=\{P_\theta:\theta\in\Theta\}$ be dominated by $\mu$, with $p_\theta\defeq \frac{dP_\theta}{d\mu}$. A statistic $T:\mathcal{X}\to\mathcal{T}$ is sufficient iff there exist measurable $g_\theta:\mathcal{T}\to[0,\infty)$ and $h:\mathcal{X}\to[0,\infty)$ such that $p_\theta(x)=g_\theta(T(x))h(x)$ $\mu$-a.e. for all $\theta$.
\begin{proof}
\begin{enumerate}[label=\roman*.]
\setlength{\itemsep}{0pt}
\item[$(\Rightarrow)$] If $T$ is sufficient, write $\nu_\theta\defeq\mathcal{L}_\theta(T)$ and choose a $\theta$-free kernel $K_t=\mathcal{L}_\theta(X\mid T=t)$,\\
so $P_\theta(A)=\int K_t(A)\nu_\theta(dt)$. Choosing densities for $\nu_\theta$ and for $K_t$ on each fiber gives $p_\theta(x)=g_\theta(T(x))h(x)$.
\item[$(\Leftarrow)$] If $p_\theta=g_\theta\circ T\,h$, then on the fiber $T=t$ the factor $g_\theta(t)$ cancels:\\
for $c(t)\defeq\int_{T^{-1}\{t\}}h\,d\mu_t>0$, $\mathbb{P}_\theta(X\in A\mid T=t)=c(t)^{-1}\int_{A\cap T^{-1}\{t\}}h\,d\mu_t$, independent of $\theta$.
\end{enumerate}
\hlt{Structural Commentary}: Factorisation separates parameter information $g_\theta(T)$ from residual shape $h$.\\
\hlt{Key Insight}: Sufficiency keeps exactly the likelihood's $\theta$-dependent part.
\end{proof}
\end{theorem}

\begin{example} \label{ex:bernoulli-normal-sufficient-statistics} \hlt{Bernoulli and Normal Sufficient Statistics}\\
Let $X_1,\ldots,X_n$ be independent and let $x=(x_1,\ldots,x_n)$.
\begin{enumerate}[label=(\roman*)]
\setlength{\itemsep}{0pt}
\item If $X_i\sim\operatorname{Bernoulli}(\theta)$, then $p_\theta(x)=\theta^s(1-\theta)^{n-s}$ with $s\defeq\sum\limits_i x_i$.\\
Hence $T=\sum\limits_iX_i$ is sufficient for $\theta$, and $\mathcal{L}_\theta(X\mid T=t)$ is uniform on $\{x:\sum\limits_i x_i=t\}$.
\item If $X_i\sim N(\mu,\sigma^2)$ with $(\mu,\sigma^2)$ unknown, then $p_{\mu,\sigma^2}(x)\propto(\sigma^2)^{-\frac{n}{2}}\exp\{-\frac{\sum\limits_i x_i^2-2\mu\sum\limits_i x_i+n\mu^2}{2\sigma^2}\}$.\\
Hence $T=(\sum\limits_iX_i,\sum\limits_iX_i^2)$ is sufficient for $(\mu,\sigma^2)$.
\end{enumerate}
\end{example}

\begin{corollary} \label{cor:mle-function-of-sufficient-statistic} \hlt{Likelihood Estimators Use Sufficient Statistics}\\
Let $\mathcal{P}=\{P_\theta:\theta\in\Theta\}$ be dominated with densities $p_\theta$, and let $T$ be sufficient with $p_\theta(x)=g_\theta(T(x))h(x)$.\\
If $a(t)\in\arg\max\limits_{\theta\in\Theta}g_\theta(t)$ exists, then $\widehat{\theta}(X)\defeq a(T(X))$ is an MLE.
\begin{proof}
For fixed $x$, $L_x(\theta)=p_\theta(x)=h(x)g_\theta(T(x))$. If $h(x)>0$, then $\arg\max\limits_\theta L_x(\theta)=\arg\max\limits_\theta g_\theta(T(x))$.\\
If $h(x)=0$, then $L_x(\theta)=0$ for all $\theta$. Thus an MLE may be chosen as $a(T(x))$.\\
\hlt{Structural Commentary}: The $\theta$-free factor $h(x)$ is irrelevant to maximization.\\
\hlt{Key Insight}: MLEs can depend on data only through a sufficient statistic.
\end{proof}
\end{corollary}

\begin{theorem} \label{thm:rao-blackwell-improvement} \hlt{Rao--Blackwell Improvement}\\
Let $T$ be sufficient and let $\widehat{\psi}\in L^2(P_\theta)$ estimate a real target $\psi(\theta)$. Choose the $\theta$-free version $\widetilde{\psi}(T)\defeq\mathbb{E}_\theta(\widehat{\psi}\mid T)$.\\
Then $\mathbb{E}_\theta\widetilde{\psi}=\mathbb{E}_\theta\widehat{\psi}$ and $\mathbb{E}_\theta(\widetilde{\psi}-\psi(\theta))^2\leq\mathbb{E}_\theta(\widehat{\psi}-\psi(\theta))^2$ for all $\theta$, with equality iff $\widehat{\psi}=\widetilde{\psi}$ $P_\theta$-a.s.
\begin{proof}
Sufficiency makes $\mathcal{L}_\theta(X\mid T)$ independent of $\theta$, so $\widetilde{\psi}$ is a $T$-statistic. Also $\mathbb{E}_\theta\widetilde{\psi}=\mathbb{E}_\theta\widehat{\psi}$.\\
Since $\mathbb{E}_\theta(\widehat{\psi}-\widetilde{\psi}\mid T)=0$ and $\widetilde{\psi}-\psi(\theta)$ is $T$-measurable, $\mathbb{E}_\theta[(\widehat{\psi}-\widetilde{\psi})(\widetilde{\psi}-\psi(\theta))]=0$.\\
Hence $\mathbb{E}_\theta(\widehat{\psi}-\psi(\theta))^2=\mathbb{E}_\theta(\widehat{\psi}-\widetilde{\psi})^2+\mathbb{E}_\theta(\widetilde{\psi}-\psi(\theta))^2\geq\mathbb{E}_\theta(\widetilde{\psi}-\psi(\theta))^2$, with equality iff $\mathbb{E}_\theta(\widehat{\psi}-\widetilde{\psi})^2=0$.\\
\hlt{Structural Commentary}: Rao--Blackwellization projects an estimator onto the $\sigma(T)$-measurable statistics.\\
\hlt{Key Insight}: Conditioning on a sufficient statistic preserves bias and removes residual noise.
\end{proof}
\end{theorem}

\begin{definition} \label{def:exponential-family-natural-statistic} \hlt{Exponential Family}
\begin{enumerate}[label=(\roman*)]
\setlength{\itemsep}{0pt}
\item A dominated model $\{P_\eta:\eta\in\mathcal{N}\}$ is a \hlt{$k$-parameter exponential family}\\
if $\mathcal{N}\subseteq\R^k$ and $p_\eta(x)=\exp\{\eta^\top T(x)-A(\eta)\}h(x)$.
\item $\eta$ is the \hlt{natural parameter}, $T(X)$ is the \hlt{natural statistic}, and $h$ is the carrier density.
\item $A(\eta)$ is the \hlt{log-partition function}: $A(\eta)=\log\int\exp\{\eta^\top T(x)\}h(x)\,d\mu(x)$ when finite.
\end{enumerate}
\end{definition}

\begin{proposition} \label{prop:exponential-family-moments} \hlt{Moment Identities in Full Exponential Families}\\
Let $p_\eta(x)=\exp\{\eta^\top T(x)-A(\eta)\}h(x)$ be a full exponential family with open natural parameter space, and assume differentiation under the integral is valid. Then $\nabla A(\eta)=\mathbb{E}_\eta T(X)$ and $\nabla^2A(\eta)=\operatorname{Cov}_\eta(T(X))$.
\begin{proof} From $1=\int p_\eta\,d\mu$, differentiation gives $0=\int(T-\nabla A(\eta))p_\eta\,d\mu$, hence $\nabla A(\eta)=\mathbb{E}_\eta T(X)$.\\
For coordinates $i,j$, differentiating $\mathbb{E}_\eta T_i=\int T_ip_\eta\,d\mu$ gives $\partial_j\mathbb{E}_\eta T_i=\mathbb{E}_\eta[T_i(T_j-\partial_jA(\eta))]=\operatorname{Cov}_\eta(T_i,T_j)$, QED.\\
\hlt{Structural Commentary}: The log-partition function cumulatively encodes the natural statistic.\\
\hlt{Key Insight}: First and second derivatives of $A$ give mean and covariance.
\end{proof}
\end{proposition}

\begin{definition} \label{def:local-asymptotic-normality} \hlt{Local Asymptotic Normality}\\
For experiments $\{P_\theta^{(n)}:\theta\in\Theta\subseteq\R^d\}$, set $\theta_{n,h}\defeq\theta_0+\frac{h}{\sqrt{n}}$.
\begin{enumerate}[label=(\roman*)]
\setlength{\itemsep}{0pt}
\item The sequence is \hlt{locally asymptotically normal} at $\theta_0$ with information $I(\theta_0)$ if,\\
for each fixed $h\in\R^d$, $\log(\frac{dP_{\theta_{n,h}}^{(n)}}{dP_{\theta_0}^{(n)}})=h^\top\Delta_n-\frac{1}{2}h^\top I(\theta_0)h+o_{P_{\theta_0}^{(n)}}(1)$.
\item The \hlt{central sequence} satisfies $\Delta_n\Rightarrow N(0,I(\theta_0))$ under $P_{\theta_0}^{(n)}$.
\end{enumerate}
\end{definition}

\begin{remark} \label{rem:lan-bridge-to-estimation-tests} \hlt{Why Local Asymptotics Matter}\\
LAN replaces local likelihood ratios by the Gaussian quadratic experiment $h^\top\Delta-\frac{1}{2}h^\top I(\theta_0)h$, giving normal approximations for estimators, tests, and confidence regions.
\end{remark}
